\section{第六天}
\begin{enumerate}
  \item 全面依法治国的重要保障（法律人和少数人）。
  \item 核心要义之\bs{坚持建设德才兼备的高素质法治工作队伍}。
        \begin{enumerate}
          \item 概念和价值。
                \begin{enumerate}
                  \item 建队伍是全推国的一项\bs{长期基础性工作}。
                  \item \bs{十年树木，百年树人}，全治国的各项工作，\bs{离不开}法工队的参与。
                  \item 高工队是我们深入贯彻党的领导，
                        顺利推进立执司守带修饰工作、打造严密法治监督体系的\bs{组织和人才保障}。
                  \item 必须要\zx{立足长远、科学规划，不断创新}法治人才培养模式，
                  \item \zd{提高}法工队的\zx{思想政治素质、业务工作能力、职业道德水准}，
                  \item 着力建设\bs{一支四忠}的社法工队，\bs{忠于党、忠于国家、忠于人民、忠于法律}。
                  \item 实现法工队的\zx{革命、正规、专职化}，\bs{革命化、正规化、专业化、职业化}。
                  \item \zd{培养造就}一大批\zx{高法人及后备力量}。
                \end{enumerate}
          \item 如何（主要背第一句，解释部分可以自己发挥）。
                \begin{enumerate}
                  \item 要加强法治\bs{专门工作队伍}建设。
                        \begin{enumerate}
                          \item 全方位\zx{提升}立法工作者、法官、检察官等公职人员的思想政治素质、业务工作能力和职业道德水准，
                          \item 强化\zx{廉政教育}，杜绝办\zx{“金钱案”、“权力案”、“人情案”}。
                        \end{enumerate}
                  \item 要加强\bs{法律服务队伍}建设。
                        \begin{enumerate}
                          \item 引导法律服务工作者\zx{坚持正确政治方向，依法依规诚信执业}，
                          \item 建立激励法律服务人才\zx{跨区域流动}机制，
                          \item 逐步解决\zx{基层和欠发达}地区法律服务\zx{资源不足和人才匮乏}等问题。
                        \end{enumerate}
                  \item 要创新\bs{法律人才培养}机制。
                        \begin{enumerate}
                          \item 坚持以法治思想为指导，\zx{立德树人，德法兼修}，培养大批高法人，
                          \item 重视\zx{高校}的法律人才培养\zx{第一阵}地作用，
                                加强\zx{法学学科体系建设}，
                                推进不同部门间的\zx{法律人才交流}，
                                全方位提升\zx{法律人才培养能力}。
                        \end{enumerate}
                \end{enumerate}
        \end{enumerate}
  \item 核心要义之\bs{坚持抓住领导干部这个“关键少数”}。
        \begin{enumerate}
          \item 坚持抓住领导干部这个“关键少数”，是全推国的\bs{关键问题}。
          \item 领导干部具体\bs{行使}党的执政权和国家的立法权、行政权、监察权、司法\bs{权}，
          \item 他们既可以对全治国的各项工作产生积极的\zx{推动}作用，
                也可以对全治国产生致命的\zx{破坏}作用，
          \item 领导干部\bs{是}全治国的\bs{关键}。
          \item 全推国\bs{必须抓住}领导干部这个“关键少数”，
                \begin{enumerate}
                  \item 不断提高他们运用\zx{法治思维和法治方式}深化改革、推动发展、化解矛盾、维护稳定的能力，
                  \item 牢固树立\bs{法定职责必须为、法无授权不可为}的法治\zx{观念}，要求他们做\zx{尊法、学法、守法、用法}的模范。
                  \item 要坚持依法\zx{治权}，用宪法和法律法规\zx{设定权力、规范权力、制约权力、监督权力}，\bs{把权力关进法律和制度的笼子里}，
                  \item 确保\zx{赋权在轨谋福、利}。\bs{党和人民赋予的权利，始终在法轨上为人民谋幸福、谋福利}。
                \end{enumerate}
        \end{enumerate}
\end{enumerate}